
Dataset reproducibility depends on documentation completeness, yet no quantitative study has linked repository UX design to metadata outcomes at scale. This work introduces friction score (0-4: automated extraction + templates + validation + API) as a quantifiable metric capturing repository design quality. Through cross-platform analysis of 10,000 datasets across OpenML, HuggingFace Datasets, and UCI ML Repository, we demonstrate that friction-reduction features correlate with metadata presence. HuggingFace (friction=3) shows 55-66\% optional field presence versus UCI (friction=0) showing 0-15\%, with differences of 45-61 percentage points (preprocessing\_code 61pp, data\_source\_url 51pp, collection\_date 45pp; p<0.0001, Cohen's h 1.0-1.8). Required fields (license, version) show 75-95\% presence across platforms with weaker friction effects (12-20pp differences, Cramér's V 0.1-0.2), validating enforcement as a distinct but coupled mechanism. This work provides the first quantitative evidence linking friction reduction to metadata completeness at 10,000+ dataset scale.
